%----------------------------------------------------------
% Thato Baruti - Master CV for portfolio
% Imperial-style two-page black-and-white CV
%----------------------------------------------------------
\documentclass[10pt,a4paper]{article}

\input{glyphtounicode}
\pdfgentounicode=1
\usepackage[T1]{fontenc}
\usepackage[utf8]{inputenc}
\usepackage{helvet}
\renewcommand{\familydefault}{\sfdefault}
\usepackage[a4paper,top=0.92cm,bottom=0.86cm,left=1.16cm,right=1.16cm]{geometry}
\usepackage{microtype}
\usepackage{parskip}
\setlength{\parskip}{2.0pt}
\usepackage{titlesec}
\usepackage{enumitem}
\usepackage{hyperref}
\hypersetup{colorlinks=true,urlcolor=black,linkcolor=black}
\pagestyle{empty}
\setlist[itemize]{leftmargin=1.08em,itemsep=1.05pt,topsep=1pt,parsep=0pt,partopsep=0pt,label={\scriptsize\textbullet}}
\titleformat{\section}{\bfseries\small}{}{0em}{\MakeUppercase}[\vspace{-1pt}\rule{\linewidth}{0.42pt}]
\titlespacing{\section}{0pt}{5.5pt}{2.0pt}
\sloppy

\begin{document}
\fontsize{9.65}{10.95}\selectfont

\begin{center}
{\fontsize{21.5}{23.5}\selectfont\textbf{THATO BARUTI}}\par
\footnotesize +44 07759 037559 $\mid$ \href{mailto:tbofficial0508@gmail.com}{tbofficial0508@gmail.com} $\mid$ Cranfield, United Kingdom $\mid$
\href{https://www.linkedin.com/in/thato-baruti-63b290222/}{LinkedIn} $\mid$
\href{https://tbofficial0508-jpg.github.io/ThatoBaruti619.github.io-/}{Portfolio} $\mid$
\href{https://github.com/tbofficial0508-jpg}{GitHub} $\mid$
\href{https://github.com/tbofficial0508-jpg/cluster-count-ai4s}{ClusterCount-AI4S}
\end{center}
\vspace{-5pt}

\section{Profile}
MSc Computational Fluid Dynamics student specialising in AI-for-engineering, CFD, scientific software and design optimisation. I build validated models and tools that connect physics, data, geometry, uncertainty and user-facing deployment, with a strong interest in roles across scientific ML, simulation, AI consulting, aerospace design and computational engineering.

\section{Education}
\textbf{Cranfield University} \textbf{(Distinction expected)} \hfill Sep 2025--Jun 2027 expected\par
\textbf{MSc Computational Fluid Dynamics}. High-scoring modules include Data Analysis and Post-Processing 87\%, Fluid Mechanics and Heat Transfer 87\%, Numerical Modelling for Incompressible Flows 84\%, plus Grid Generation, Numerical Methods, Turbulence Modelling and High Performance Computing.

\textbf{University of Sheffield} \textbf{(2nd Class Honours)} \hfill Sep 2022--Jun 2025\par
\textbf{BEng Aerospace Engineering}. Relevant study in Aerodynamic Design, CFD, Aircraft Dynamics and Control, Propulsion, Aerospace Design, CAD, MATLAB, project management and technical communication. Awarded the Botswana Government Top Achievers Scholarship and Sheffield International Undergraduate Scholarship.

\section{Scientific AI, CFD and Software Projects}
\textbf{AeroGenic Platform and PiSA Surrogate Algorithm} \hfill Cranfield University thesis, 2025--ongoing
\begin{itemize}
  \item Developed the engineering-tool deployment layer of my thesis Physics-informed Surrogate Algorithm for inverse-design optimisation of 3D UAV wings, connecting parametric geometry, aerodynamic data, PyTorch graph-style modelling, optimisation, validation and API/UI deployment.
  \item Trained on a 300-case panel design space and anchored to RANS validation; complete GPU inference returns aerodynamic coefficients, uncertainty, surface pressure and decoded section geometry in 0.16 ms, about six orders faster than the matching 64-core RANS workflow.
  \item Improved drag/lift accuracy by 76\% over a panel-method baseline and recovered held-out RANS surface pressure to 7.6\%, while documenting validation gates, residual-correction limits, uncertainty coverage and behaviour under unseen airfoil-family shift.
  \item Built a FastAPI/Gradio tool for requirements, geometry import/export, prediction, confidence reporting and optimisation, translating research models into a usable engineering workflow for stakeholder review.
\end{itemize}

\textbf{ThermoPulseLBM Scientific Solver} \hfill Personal C++/HPC project, 2026
\begin{itemize}
  \item Built a C++20 Lattice Boltzmann solver with D2Q9 flow, D2Q5 thermal/passive-scalar transport, serial, OpenMP and OpenCL backends, CMake packaging, benchmark scripts and 67 local tests passing.
  \item Validated Poiseuille, lid-driven cavity, Womersley and Graetz benchmarks, including Poiseuille L2 error 4.9e-5, Ghia cavity deviations of 0.52--0.82\%, Womersley amplitude error 0.11\% and Graetz Nusselt error 0.23\%.
\end{itemize}

\textbf{ClusterCount-AI4S Public GitHub Project} \hfill Personal AI-for-science project, 2026
\begin{itemize}
  \item Built a validation-first microscopy image counting framework using Python, scikit-learn, synthetic controls, interpretable extremely randomized trees, overlays, model comparison and warning signals.
  \item Achieved 4.67\% MAPE on 300 held-out synthetic microscopy images and 4.45\% MAPE inside a moderate blur/clustering envelope; reported operating-range limits, disagreement signals and external-data failure modes rather than hiding deployment risk.
\end{itemize}

\textbf{OpenFOAM RANS Validation and CFD Workflow Automation} \hfill Cranfield University, 2025--ongoing
\begin{itemize}
  \item Built OpenFOAM RANS validation cases for selected aerodynamic designs, applying mesh-quality checks, $y^+$ control, coefficient stability, force extraction, residual diagnostics and Grid Convergence Index reasoning.
  \item Automated case preparation and post-processing support across multiple CFD runs, producing engineering reports, publication-style figures and supervisor/stakeholder updates.
\end{itemize}

\newpage

\section{Engineering Design and Team Projects}
\textbf{Design Optimisation of Pod Deployable UAV} \hfill Cranfield University, 2026
\begin{itemize}
  \item Contributed to an assessed group design project awarded 82\%, with 85\% for the final presentation; coordinated task management in Notion and translated mission requirements into technical workstreams.
  \item Developed a Physics-Informed DGCNN FiLM-SIREN surrogate for EDF-body aerodynamic prediction to single-digit error margins below 8.2\%, identifying initialization dilution as a failure mode and implementing a two-phase training protocol with spatially weighted physics regularisation.
  \item Produced a preliminary SolidWorks design that became the backbone of the EDF fan model and supported high-fidelity UCNS3D simulations, balancing design intent, numerical evidence and presentation clarity.
\end{itemize}

\textbf{AeroMAP UAS Build} \hfill Performance Lead, University of Sheffield, Oct 2023--May 2024
\begin{itemize}
  \item Led simulation-led performance improvement across five aerodynamic and structural components using CAD iteration, ANSYS CFD, FEA and physical flight-test feedback.
  \item Reported 25\% greater stability and 15\% greater endurance, balancing aerodynamic performance, structural practicality, manufacturability and testing evidence before recommending changes.
\end{itemize}

\textbf{SunStratos High-Altitude Surveillance Glider} \hfill Aerostructures Lead, Oct 2023--Jul 2025
\begin{itemize}
  \item Led 12 students through airframe concept development, CAD modelling, structural integrity work, manufacturing planning and prototype delivery, improving delivery speed by 33\%.
  \item Converted broad mission aims into component responsibilities, design reviews and build tasks, strengthening communication across aerodynamics, structures and manufacturing.
\end{itemize}

\textbf{UKSEDS National Rocketry Competition} \hfill Design Team Lead, Oct 2022--May 2023
\begin{itemize}
  \item Owned the CFD-led design workflow from CAD geometry through meshing, solver setup, post-processing and technical reporting; reduced drag coefficient by 22\% and ranked Top 10 nationally among 30+ teams.
\end{itemize}

\section{Programmes, Leadership and Additional Experience}
\textbf{Orange Summer Challenge, Selected Young Talent - AI/ML Track} \hfill Jul 2026--present
\begin{itemize}
  \item Selected for Orange Digital Center Botswana's AI as a Business Accelerator programme, focused on applying AI/ML to real business challenges across fintech, energy, customer experience and generative AI.
\end{itemize}

\textbf{Royal Aeronautical Society Student Affiliate} \hfill Oct 2025--present
\begin{itemize}
  \item Build professional aerospace awareness through technical events, networking and continued learning across design, simulation and aircraft development.
\end{itemize}

\textbf{Sheffield Engineering Leadership Academy Cohort} \hfill Sep 2023--Jun 2025
\begin{itemize}
  \item Developed leadership, project delivery and communication skills through a selective engineering leadership programme and multidisciplinary team activities.
\end{itemize}

\textbf{Formula 1 Silverstone Steward, Sierra 1 Security Stewarding Ltd} \hfill Jun 2023
\begin{itemize}
  \item Used calm communication, radio coordination and judgement under pressure to guide spectators, resolve high-pressure queries and escalate safety issues at a major Formula 1 event.
\end{itemize}

\section{Professional Development}
\textbf{Technical training and portfolio preparation} \hfill 2023--2026
\begin{itemize}
  \item Completed Introduction to Tecplot and pyTecplot plus MATLAB Onramp; use Tecplot, ParaView, MATLAB and Python to convert simulation output into clear engineering evidence, plots and technical presentations.
  \item Completed Cranfield climate and energy-transition study, strengthening interest in roles where simulation, AI and engineering judgement support net-zero design, renewable energy and resilience decisions.
  \item Maintain public portfolio and GitHub evidence across AeroGenic, ThermoPulseLBM and ClusterCount, including project READMEs, validation notes, screenshots, reports and reproducible scripts where confidentiality allows.
\end{itemize}
\section{Technical Skills}
\textbf{AI and data science:} PyTorch, PyTorch Geometric-style graph ML, scikit-learn, pandas, NumPy, model evaluation, supervised learning, uncertainty checks, residual correction, optimisation, validation dashboards, failure-slice diagnostics and responsible model-limit reporting.\par
\textbf{Scientific software and deployment:} Python, C++20, MATLAB, Bash, Git/GitHub, FastAPI, Gradio, CMake, OpenMP, OpenCL, Linux/HPC workflows, Excel, PowerPoint, technical reports, demo outputs and reproducible benchmark scripts.\par
\textbf{Simulation and design:} OpenFOAM, ANSYS Fluent/CFX exposure, UCNS3D, FluidX3D, VSPAERO/panel methods, finite-volume methods, LBM, RANS, heat transfer, turbulence modelling, mesh-quality checks, GCI, $y^+$, ParaView, Tecplot, SolidWorks, Fusion 360 and CAD geometry preparation.

\section{Portfolio Themes}
Current portfolio evidence spans physics-informed surrogate modelling, CFD validation, solver development, AI-for-science evaluation, UAV design optimisation, aerodynamic performance, technical communication and client-style demos. I am particularly interested in work where modern AI can accelerate engineering decisions while remaining transparent, validated and physically credible.

\end{document}
